\documentclass[10pt,a4paper]{amsart}
\usepackage[margin=19mm]{geometry}
\usepackage{amsmath,amssymb,mathtools}
\usepackage[scr=rsfs,cal=euler]{mathalfa}
\usepackage{xfrac}
\usepackage[unicode,hidelinks,bookmarks=false]{hyperref}
\newcommand{\ie}{i.e.\ }
\newcommand{\cf}{cf.\ }
\newcommand{\resp}{respectively\ }
\newcommand{\Subsection}{\subsection}
\providecommand{\nobreakdash}{\nobreak\hbox{-}}
\setlength{\emergencystretch}{2em}
\multlinegap0pt
\usepackage{amssymb}
\usepackage{bm}
\usepackage{enumerate}
\newtheorem{lemma}{Lemma}
\newcommand{\cl}{\mathcal K}
\title{Large Cliques and Clique Spectral Radius in the Erd\H{o}s--S\'{o}s Problem}
\author{XiaoJun Zhao, YueJian Peng}
\date{Source: arXiv:2608.25746v1 (CC BY 4.0)}
\begin{document}
\maketitle

\noindent\emph{Source and licence.} Large Cliques and Clique Spectral Radius in the Erd\H{o}s--S\'{o}s Problem. Version arXiv:2608.25746v1 (CC BY 4.0), distributed by arXiv under the Creative Commons Attribution 4.0 International licence, http://creativecommons.org/licenses/by/4.0/, as recorded in the arXiv licence metadata for this exact version and shown on the abstract page https://arxiv.org/abs/2608.25746v1. Full licence notice: \texttt{licenses/arxiv-2608.25746-CC-BY-4.0.txt}.\par\smallskip

\noindent\textit{Editorial scope.} Counted unit: the article's lemma component-bound (source0.tex lines 467--475). The article's complete original proof of it is delivered with it (source0.tex lines 477--583, one \texttt{proof} environment of 107 lines). No mathematics is written, repaired or re-derived: the statement and the proof are the article's own lines, in the article's own notation, and no retained line is substituted or re-flowed.  Retained uncounted as the source-local context the proof needs: lines 415--420; lines 426--434; lines 455--461.  Nothing was omitted from any retained span, so the package carries no omission record.  No retained line contains a citation, so the wrapper carries no bibliography and no bibliography entry.  No retained line contains a figure environment, an image-inclusion command or drawing code, so no image is delivered and the package declares no asset dependency.  Editorial formatting only, in addition: an AMS \texttt{amsart} preamble that restates the article's own declarations and macros used by the retained lines, namely the environments \texttt{lemma} on separate counters, together with the standard packages those lines need.  The retained environments print in this document as Lemma 1 in order of appearance, whereas the article prints the same items with the numbering of its own full document.  The complete native source of the article as distributed in the arXiv e-print for this exact version is bundled unchanged as \texttt{sources/208594/source0.tex}.
\begin{lemma}\label{large-overlap}
Let \(d\ge2\) and \(t\ge 3d-1\). Let \(T_t\) be a \(t\)-vertex tree, and let \(G\) be \(T_t\)-free. If \(K\) and \(K'\) are distinct intersecting \((t-d)\)-cliques in \(G\), then
\[
t-2d+1\le |V(K)\cap V(K')|\le t-d-1.
\]
\end{lemma}

\begin{lemma}\label{large-transitivity}
Let \(\cl\) be the family of \((t-d)\)-cliques in a \(T_t\)-free graph, where \(d\ge2\) and \(t\ge3d-1\). Then the relation
\[
R\sim R'
\quad\Longleftrightarrow\quad
V(R)\cap V(R')\ne\emptyset
\]
is an equivalence relation on \(\cl\).
\end{lemma}

\begin{lemma}\label{pg}
Let \(G\) be a graph in which every edge is contained in an \(r\)-clique. Let \(\cl\) be the set of all \(r\)-cliques in \(G\), and define a relation on \(\cl\) by \(R\sim R'\) if and only if \(V(R)\cap V(R')\ne\emptyset\). If \(\sim\) is an equivalence relation with classes \(\cl_1,\ldots,\cl_s\), then
\[
G[V(\cl_1)],\ldots,G[V(\cl_s)]
\]
are precisely the nontrivial components of \(G\). Here \(V(\cl_i)\) is the union of the vertex sets of the cliques in \(\cl_i\).
\end{lemma}

\begin{lemma}\label{component-bound}
Let $d\ge2$, let $T_t$ be a $t$-vertex tree, and put $r=t-d$. For a $T_t$-free graph $G$, let $G_r$ be the spanning subgraph of $G$ whose edges are precisely those contained in at least one copy of $K_r$. If
\[
\lambda(T_t)\le d-1
\qquad\text{and}\qquad
 t\ge d^2-d+3,
\]
then every nontrivial component of $G_r$ has order at most $t-1$.
\end{lemma}

\begin{proof}
Since
\[
(d^2-d+3)-(3d-1)=(d-2)^2\ge0,
\]
we have $t\ge3d-1$. By Lemma~\ref{large-transitivity}, intersection is therefore an equivalence relation on the family of $r$-cliques of $G_r$, and Lemma~\ref{pg} identifies the corresponding equivalence classes with the nontrivial components of $G_r$.

Suppose, for a contradiction, that a nontrivial component $C$ has at least $t$ vertices. Fix an $r$-clique $R$ in $C$ and choose distinct vertices
\[
u_1,\ldots,u_d\in V(C)\setminus V(R).
\]
Every vertex of $C$ lies in an $r$-clique from the equivalence class corresponding to $C$. Hence, for each $i\in[d]$, there is an $r$-clique $R_i$ containing $u_i$ and intersecting $R$. By Lemma~\ref{large-overlap},
\begin{equation}\label{component-Ai}
A_i:=N_R(u_i),
\qquad
|A_i|\ge |V(R_i)\cap V(R)|\ge t-2d+1=r-d+1.
\end{equation}
We shall embed $T_t$ into $G[V(R)\cup\{u_1,\ldots,u_d\}]$.

Let $\ell$ be the number of leaves of $T_t$.

\smallskip
\noindent\emph{Case 1: $\ell\ge d$.}

Choose any $d$ leaves, denote them by $v_1,\ldots,v_d$, map $v_i$ to $u_i$, and let $P$ be the set of their distinct parents. Put $p=|P|$. Since $\lambda(T_t)\le d-1$, we have $p\ge2$. For $y\in P$, define
\[
I(y)=\{i\in[d]:y\text{ is the parent of }v_i\},
\qquad
B_y=\bigcap_{i\in I(y)}A_i.
\]
Since the positive integers $|I(y)|$, $y\in P$, sum to $d$, we have $|I(y)|\le d-p+1$. By~\eqref{component-Ai},
\[
|B_y|\ge r-|I(y)|(d-1)
\ge t-d-(d-p+1)(d-1).
\]
Moreover,
\[
\begin{aligned}
|B_y|-p
&\ge t-d^2-d+1+p(d-2).
\end{aligned}
\]
Since $p\ge2$ and $t\ge d^2-d+3$, the right-hand side is at least
\[
t-d^2+d-3\ge0.
\]
Hence $|B_y|\ge p$ for every $y\in P$. Thus,  there are distinct vertices $b_y\in B_y$. Map each parent $y$ to $b_y$ and extend the map to a bijection from $V(T_t)\setminus\{v_1,\ldots,v_d\}$ onto $V(R)$. All tree edges are preserved, yielding a copy of $T_t$, a contradiction.

\smallskip
\noindent\emph{Case 2: $\ell<d$.}

Since $l\ge 2$, we have $d\ge3$. Let $b$ be the number of vertices of degree at least $3$ and let $h$ be the number of vertices whose degree is not $2$. The identity
\[
\ell=2+\sum_{v:\,d_{T_t}(v)\ge3}\bigl(d_{T_t}(v)-2\bigr)
\]
implies $b\le\ell-2$, and hence
\[
h=\ell+b\le2\ell-2.
\]
Suppress all degree-$2$ vertices. The resulting skeleton is a tree on $h$ vertices, whose $h-1$ edges correspond to the maximal bare paths of $T_t$.

On every maximal bare path, mark the degree-$2$ vertex closest to each end whenever it exists; if an end is a leaf, also mark the second degree-$2$ vertex from that end whenever it exists. Thus at most $2(h-1)+\ell$ degree-$2$ vertices are marked. If $M$ denotes the number of unmarked degree-$2$ vertices, then
\[
M\ge t-h-2(h-1)-\ell=t-3h+2-\ell\ge t-7\ell+8.
\]
From each remaining path choose every third unmarked vertex. The marking at the ends guarantees that choices from distinct bare paths are also at distance at least $3$, and no chosen vertex is within distance $2$ of a leaf. Moreover,
\[
\begin{aligned}
M-[3(d-\ell)-2]
&\ge t-3d-4\ell+10\\
&\ge(d^2-d+3)-3d-4(d-1)+10\\
&=d^2-8d+17=(d-4)^2+1>0.
\end{aligned}
\]
Thus we can choose $d-\ell$ degree-$2$ vertices $x_1,\ldots,x_{d-\ell}$ which are pairwise at distance at least $3$, and whose neighbors are all distinct and are different from all parents of leaves.

Map all $\ell$ leaves and the vertices $x_1,\ldots,x_{d-\ell}$ bijectively to $u_1,\ldots,u_d$. Re-indexing if necessary, suppose the leaves are mapped to $u_1,\ldots,u_\ell$. Let $P$ be the set of distinct parents of the leaves and put $p=|P|$. For $y\in P$, define
\[
I(y)=\{i\in[\ell]:y\text{ is the parent of the leaf mapped to }u_i\},
\qquad B_y=\bigcap_{i\in I(y)}A_i.
\]
Since $|I(y)|\le\ell-p+1$ and $\ell\le d-1$,
\[
\begin{aligned}
|B_y|-p
&\ge t-d-(\ell-p+1)(d-1)-p\\
&\ge t-d-(d-p)(d-1)-p\\
&\ge p(d-2)-d+3\ge0.
\end{aligned}
\]
For each selected degree-$2$ vertex mapped to $u_j$, prescribe the candidate set $A_j$ for each of its two neighbors. The total number of prescribed neighbor-images is
\[
p+2(d-\ell)\le2d-\ell\le2d-2.
\]
On the other hand, by~\eqref{component-Ai},
\[
|A_j|\ge t-2d+1\ge2d-2,
\]
because
\[
t-2d+1-(2d-2)\ge d^2-5d+6=(d-2)(d-3)\ge0.
\]
%Hall's condition follows. Indeed, a subfamily containing one of the sets $A_j$ has union of size at least $2d-2$, which is at least the total number of prescribed neighbor-images; a subfamily consisting only of the sets $B_y$ has union of size at least $p$, and there are only $p$ such sets.
Therefore all parents of leaves and all neighbors of the selected degree-$2$ vertices can be mapped to distinct vertices of $R$ in their prescribed candidate sets. Extend this map to a bijection from the remaining $t-d$ vertices of $T_t$ onto $V(R)$. Again all tree edges are preserved, giving a copy of $T_t$, a contradiction.

Both cases are impossible, so every nontrivial component of $G_r$ has order at most $t-1$.
\end{proof}

\end{document}
